\begin{tikzpicture}[
  box/.style={draw, rounded corners, thick, align=center, minimum width=32mm, minimum height=10mm, fill=black!3},
  note/.style={draw, rounded corners, align=left, inner sep=4pt, fill=black!2},
  labelnote/.style={align=center, fill=white, inner sep=1pt, font=\small},
  arrow/.style={-{Latex[length=2.5mm]}, thick, shorten <=2pt, shorten >=2pt}
]
\node[box] (access) at (0,0) {State access\\and scheduling};
\node[box] (exec) at (4.3,0) {State-aware\\execution};
\node[box] (evo) at (8.6,0) {State evolution\\and reuse};
\node[labelnote] (l1) at (2.15,0.8) {conflicts, locality,\\queueing};
\node[labelnote] (l2) at (6.45,0.8) {cost, quality,\\resource usage};
\node[labelnote] (fb) at (4.3,2.0) {future hotspots,\\memory pressure};
\node[note, text width=108mm] (note1) at (4.3,-2.15) {
\textbf{Propagation view.} Local mismatches rarely stay local. Hidden access contention distorts execution decisions; poorly modeled execution reshapes what can be updated or retained; aggressive updates feed back into later access pressure and reuse quality.
};

\coordinate (topL) at (-0.4,1.45);
\coordinate (topR) at (9.0,1.45);

\draw[arrow] (access.east) -- (exec.west);
\draw[arrow] (exec.east) -- (evo.west);
\draw[arrow] (evo.north) |- (topR) -- (topL) |- (access.north);
\draw[arrow] (exec.south) -- (note1.north);
\end{tikzpicture}
